%\documentclass[10pt]{article}
\documentclass[tikz,border=3pt]{standalone}
%\usepackage[usenames]{color} %used for font color
\usepackage{amssymb} %maths
\usepackage{amsmath} %maths
\usepackage[utf8]{inputenc} %useful to type directly diacritic characters
%\usepackage[svgnames]{xcolor}
\usepackage{tikz}
\usetikzlibrary{decorations.markings}
\pagestyle{empty}

\usepackage{amsmath, amssymb, amsfonts,amsthm} % Math packages for equations
\input{/Users/deepak/ownCloud/root/research/articles/timesarrow/bridgeman-chubbs/testing/macros}


\pgfdeclarelayer{edgelayer}
\pgfdeclarelayer{nodelayer}
\pgfsetlayers{edgelayer,nodelayer,main}

\tikzstyle{none}=[inner sep=2pt]

\usepackage[graphics,tight page,active]{preview}
\PreviewEnvironment{tikzpicture}
\newlength{\imagewidth}
\newlength{\imagescale}

\tikzstyle{tensor}=[rectangle,draw=blue!50,fill=blue!20,thick]

\tikzset{node distance=2cm}
\begin{document}
\definecolor{darkorange}{HTML}{f9a262}
\definecolor{lightorange}{HTML}{f2dbde}

\tikzstyle{tensor}=[rectangle,draw=blue!50,fill=blue!20,thick]

\tikzstyle{spin}=[circle,draw=black!50!white,fill=darkorange,thick]
\tikzstyle{vertex}=[circle,radius=3.0,draw=black,fill=lightorange,thick]

\tikzstyle{siteline} = [dashed,line width=0.1cm];
\tikzstyle{vertexline} = [line width=0.05cm]

\tikzset{
  spinsite/.pic = {
      \node[vertex, scale=16] (v1) at (0,0) {};

  \node[spin, scale=3] (s1) at (-1,1) {};
  \node[spin, scale=3] (s2) at (1,1) {};
  \node[spin, scale=3] (s3) at (1,-1) {};
  \node[spin, scale=3] (s4) at (-1,-1) {};

  \node[font=\huge] at ($ (s1) + (-0.6,0.6) $) {$1$};
  \node[font=\huge] at ($ (s2) + (0.6,0.6) $) {$2$};
  \node[font=\huge] at ($ (s3) + (0.6,-0.6) $) {$3$};
  \node[font=\huge] at ($ (s4) + (-0.6,-0.6) $) {$4$};

  \draw[siteline] (s1) -- node [midway, above] {$CZ$} (s2) -- node [midway,right]  {$CZ$} (s3) -- node [midway, below] {$CZ$} (s4) -- node [midway, left] {$CZ$} (s1);

  \draw[vertexline] (s1) -- ($ (s1) + (-3,0) $);
  \draw[vertexline] (s1) -- ($ (s1) + (0,3) $);
  \draw[vertexline] (s2) -- ($ (s2) + (0,3) $);
  \draw[vertexline] (s2) -- ($ (s2) + (3,0) $);
  \draw[vertexline] (s3) -- ($ (s3) + (0,-3) $);
  \draw[vertexline] (s3) -- ($ (s3) + (3,0) $);
  \draw[vertexline] (s4) -- ($ (s4) + (0,-3) $);
  \draw[vertexline] (s4) -- ($ (s4) + (-3,0) $);
  }
}

\tikzset{
    three-vertex/.pic =
        {
            \node [vertex] (v0) at (0,0) {#1};
            \draw [vertexline] (v0) -- (0,-1);
            \draw [rotate=120,vertexline] (v0) -- (0,-1);
            \draw [rotate=240,vertexline] (v0) -- (0,-1);
          }
}

\begin{tikzpicture}

  \draw [help lines] (-3,-3) grid (3,3);

  \pic at (0,0) {three-vertex={$v_0$}};



%  \pic at (0,0) {spinsite};


%  \pic at (0,0) {spinsite};
%  \pic at (0,6) {spinsite};
%  \pic at (6,0) {spinsite};
%  \pic at (6,6) {spinsite};

%  \foreach \i in {1,..,4}{
%    \draw[vertexline] (s\i) -- ($ (s\i) + (0,3) $);
%  };

\end{tikzpicture}




\end{document}